\chapter{La rete di sicurezza e il verdetto finale}\label{cap:p1-sicurezza}

% Capitolo della Parte I --- safety gate, incertezza, fallback e flusso completo.

La rete ha imparato a giudicare i pazienti, ma \textbf{può sbagliare}. In
medicina un errore non è un disturbo: mandare a casa un paziente in crisi
può essere pericoloso. Per questo il sistema non si fida ciecamente della rete,
ma le mette accanto una \textbf{rete di sicurezza}.

\section{L'ispettore di sicurezza (safety gate)}\label{sec:p1-sic-gate}

Prima che la rete dica la sua, interviene un \textbf{controllo automatico e
rigido} basato sulle regole dell'esperto. Il suo compito è semplice e
inflessibile:

\begin{itemize}
  \item se le regole dicono che il caso è \textbf{critico} (rischio alto),
        allora la risposta è \textbf{alto} e basta;
  \item in quel caso la rete \textbf{non viene nemmeno consultata}, e il
        verdetto non può essere abbassato.
\end{itemize}

In pratica: sulle cose gravi, l'ultima parola spetta sempre alle regole
dell'esperto, mai alla macchina che ha imparato per imitazione.

\section{Il pulsante «non sono sicuro» (fallback)}\label{sec:p1-sic-fallback}

Anche nei casi non critici, la rete a volte ha dei dubbi. Ogni sua risposta
arriva con un numero di \textbf{sicurezza}: quanto si fida di sé. Se la
sicurezza scende sotto una soglia (nel progetto, 0,6), la rete ammette
«non sono sicuro» e si fa dare la risposta dalle regole dell'esperto. Meglio
chiedere, che sbagliare.

\section{Tutto insieme: da un paziente al verdetto}\label{sec:p1-sic-flusso}

Mettiamo insieme i pezzi. Quando arrivano i sei numeri di un paziente, il
sistema procede così:

\begin{enumerate}
  \item L'ispettore di sicurezza guarda le regole: il caso è critico? Se sì
        → risposta \textbf{alto}, e si ferma qui (la rete non viene
        consultata).
  \item Se non è critico, tocca alla rete: lei guarda i numeri e dà un rischio
        con la sua sicurezza.
  \item La sicurezza è bassa (sotto 0,6)? Allora si chiede all'esperto
        (fallback).
  \item Altrimenti si tiene la risposta della rete.
  \item In ogni caso, il risultato viene registrato e, se il caso è grave, il
        medico viene avvisato.
\end{enumerate}

\begin{esempio}{Tre pazienti, tre finali}
\begin{itemize}
  \item Mario ha tutti i valori nella norma. L'ispettore non vede criticità,
        la rete risponde «basso» con sicurezza alta: verdetto basso.
  \item Lucia ha la pressione a 200/120. L'ispettore vede subito un valore
        critico: risposta «alto» immediata, la rete non viene nemmeno
        interpellata.
  \item Sara è al limite, la rete risponde ma con sicurezza 0,55 (sotto la
        soglia). Scatta il fallback: decide l'esperto con le sue regole.
\end{itemize}
\end{esempio}

\begin{esempio}{E alla fine?}
Ogni verdetto viene salvato in un database, insieme a chi l'ha dato (rete,
ispettore o esperto) e a quanto era sicuro. Se il caso è grave, parte una
notifica per il medico. Così niente passa inosservato.
\end{esempio}
